\subsubsection{Matrix Translation of Dirac and Pauli Component Checks}
\label{sec:matrix-qm-components}
The intrinsic fixed-projector amplitude expansion from
Eq.\ \eqref{eq:fixed-projector-amplitudes} is translated into one selected
two-entry column by
\begin{equation*}
\begin{aligned}
\pv{\psi}_{\eqtext{block}}
=a\pv{\Pi}(+)+b\pv{a}(-)
&\quad\longleftrightarrow\quad
\mat{\psi}_{\eqtext{block}}=\mcol{a,b},\\
\herm{\pv{\psi}_{\eqtext{block}}}
&\quad\longleftrightarrow\quad
\mherm{\mat{\psi}_{\eqtext{block}}}.
\end{aligned}
\end{equation*}
The projector--ladder expansion remains the intrinsic definition.

\paragraph{Pauli current.}
The large and small projected paravectors
\(\pv{\psi}_{\eqtext{Pauli\_large}}\) and
\(\pv{\psi}_{\eqtext{Pauli\_small}}\) are represented by the matrix columns
\(\mat{\psi}_{\eqtext{Pauli\_large}}\) and
\(\mat{\psi}_{\eqtext{Pauli\_small}}\), respectively. The row--column form
of the current reduction is
\begin{align*}
\vct{j}(\bPsi_{\eqtext{Dirac\_with\_potential}})
&=\mherm{\mat{\psi}_{\eqtext{Pauli\_large}}}\sigv
\mat{\psi}_{\eqtext{Pauli\_small}}\notag\\
&\quad+\mherm{\mat{\psi}_{\eqtext{Pauli\_small}}}\sigv
\mat{\psi}_{\eqtext{Pauli\_large}},\\
\mat{\psi}_{\eqtext{Pauli\_small}}
&=\frac{\sigv\cdot\vct{P}}{2m}
\mat{\psi}_{\eqtext{Pauli\_large}}+O(m^{-2}).
\end{align*}
Application of the Pauli product in the written order gives
\begin{align*}
\vct{j}(\pv{\psi}_{\eqtext{Pauli\_large}})
&=\frac{1}{2m}\Bigl(
\mherm{\mat{\psi}_{\eqtext{Pauli\_large}}}
  (\vct{P}\mat{\psi}_{\eqtext{Pauli\_large}})\notag\\
&\qquad
+\mherm{(\vct{P}\mat{\psi}_{\eqtext{Pauli\_large}})}
  \mat{\psi}_{\eqtext{Pauli\_large}}\Bigr)\notag\\
&\quad+\frac{\gradv\times
(\mherm{\mat{\psi}_{\eqtext{Pauli\_large}}}\sigv
 \mat{\psi}_{\eqtext{Pauli\_large}})}{2m}
+O(m^{-2}).
\end{align*}
In the first numerator, the three entries are obtained by replacing
\(\vct{P}\) in turn by \(P_1,P_2,P_3\). The first term is the
kinetic-momentum flow, and the second is the magnetization current. This is
the ordinary row--column translation of the intrinsic formula in
Eq.\ \eqref{eq:pauli-current}.

\paragraph{Free block Hamiltonian.}
The ordered-pair Hamiltonian from Eq.\ \eqref{eq:dirac-hamiltonian} is
packaged as
\begin{equation*}
\mat{\mathcal H}
:=
\begin{bmatrix}
-\sigv\cdot\vct{p}&m\\
m&\sigv\cdot\vct{p}
\end{bmatrix},
\qquad
\bPsi_{\eqtext{Dirac}}
\quad\longleftrightarrow\quad
\begin{bmatrix}
\bPsi_{\eqtext{Dirac},1}\\
\bPsi_{\eqtext{Dirac},2}
\end{bmatrix}.
\end{equation*}
Direct block multiplication gives
\begin{align*}
\mat{\mathcal H}^{\,2}
&=
\begin{bmatrix}
(\sigv\cdot\vct{p})^2+m^2
&-m\sigv\cdot\vct{p}+m\sigv\cdot\vct{p}\\
-m\sigv\cdot\vct{p}+m\sigv\cdot\vct{p}
&(\sigv\cdot\vct{p})^2+m^2
\end{bmatrix}\notag\\
&=(m^2+\vct{p}^2)\I
=E(\vct{p})^2\I.
\end{align*}
Because the momentum components and mass are real and each sigma generator
is Hermitian, the block operation satisfies
\begin{equation*}
\mherm{\mat{\mathcal H}}=\mat{\mathcal H}.
\end{equation*}
The energy-sign projection is consequently checked by
\begin{align*}
\mat{\mathcal H}\bPsi_{\eqtext{Dirac\_energy}}(\pm)
&=\frac12\left(
\mat{\mathcal H}\bPsi_{\eqtext{Dirac}}
\pm\frac{\mat{\mathcal H}^{\,2}\bPsi_{\eqtext{Dirac}}}
{E(\vct{p})}\right)\notag\\
&=\frac12\left(
\mat{\mathcal H}\bPsi_{\eqtext{Dirac}}
\pm E(\vct{p})\,\bPsi_{\eqtext{Dirac}}\right)\notag\\
&=\pm E(\vct{p})\,\bPsi_{\eqtext{Dirac\_energy}}(\pm),\\
\bPsi_{\eqtext{Dirac}}
&=\bPsi_{\eqtext{Dirac\_energy}}(+)\notag\\
&\quad+\bPsi_{\eqtext{Dirac\_energy}}(-).
\end{align*}
These lines are the block-array translation of
Eqs.\ \eqref{eq:dirac-hamiltonian-action} and
\eqref{eq:energy-projection-properties}.

\paragraph{Energy and momentum-axis columns.}
The positive- and negative-energy families are kept distinct. Within either
family, the two momentum-axis columns are written as
\begin{equation*}
\begin{aligned}
\bPsi_{\eqtext{Dirac\_positive\_energy\_axis},1}(\pm)
&\quad\longleftrightarrow\quad\mcol{a,b},\\
\bPsi_{\eqtext{Dirac\_positive\_energy\_axis},2}(\pm)
&\quad\longleftrightarrow\quad\mcol{c,d}.
\end{aligned}
\end{equation*}
The negative-energy family is represented in a separate calculation by
\begin{equation*}
\begin{aligned}
\bPsi_{\eqtext{Dirac\_negative\_energy\_axis},1}(\pm)
&\quad\longleftrightarrow\quad\mcol{a,b},\\
\bPsi_{\eqtext{Dirac\_negative\_energy\_axis},2}(\pm)
&\quad\longleftrightarrow\quad\mcol{c,d}.
\end{aligned}
\end{equation*}
For the formulas below, \(s=+1\) when the literal axis argument is \(+\), and
\(s=-1\) when that argument is \(-\). The momentum-axis equation is therefore
\begin{align*}
p_3a+(p_1-ip_2)b&=s\|\vct{p}\|a,
&
(p_1+ip_2)a-p_3b&=s\|\vct{p}\|b,\\
p_3c+(p_1-ip_2)d&=s\|\vct{p}\|c,
&
(p_1+ip_2)c-p_3d&=s\|\vct{p}\|d.
\end{align*}
The energy relations obtained from
Eq.\ \eqref{eq:energy-axis-block-check} are, for the positive-energy family,
\begin{align*}
mc&=\bigl(E(\vct{p})+s\|\vct{p}\|\bigr)a,
&
md&=\bigl(E(\vct{p})+s\|\vct{p}\|\bigr)b,\\
ma&=\bigl(E(\vct{p})-s\|\vct{p}\|\bigr)c,
&
mb&=\bigl(E(\vct{p})-s\|\vct{p}\|\bigr)d,
\end{align*}
and, for the negative-energy family,
\begin{align*}
mc&=\bigl(-E(\vct{p})+s\|\vct{p}\|\bigr)a,
&
md&=\bigl(-E(\vct{p})+s\|\vct{p}\|\bigr)b,\\
ma&=\bigl(-E(\vct{p})-s\|\vct{p}\|\bigr)c,
&
mb&=\bigl(-E(\vct{p})-s\|\vct{p}\|\bigr)d.
\end{align*}
For each family separately, \(a,b,c,d\) denote that family's amplitudes. The
positive-energy density and momentum-axis current obey
\begin{align*}
\rho\!\left(
\bPsi_{\eqtext{Dirac\_positive\_energy\_axis}}(\pm)
\right)
&=|a|^2+|b|^2+|c|^2+|d|^2,\\
\vct{u}\cdot\vct{j}\!\left(
\bPsi_{\eqtext{Dirac\_positive\_energy\_axis}}(\pm)
\right)
&=s\bigl(|c|^2+|d|^2-|a|^2-|b|^2\bigr).
\end{align*}
The negative-energy family has the same component form with its own
amplitudes:
\begin{align*}
\rho\!\left(
\bPsi_{\eqtext{Dirac\_negative\_energy\_axis}}(\pm)
\right)
&=|a|^2+|b|^2+|c|^2+|d|^2,\\
\vct{u}\!\cdot\!\vct{j}\!\left(
\bPsi_{\eqtext{Dirac\_negative\_energy\_axis}}(\pm)
\right)
&=s\bigl(|c|^2+|d|^2-|a|^2-|b|^2\bigr).
\end{align*}
Their normalized currents distinguish the two energy families:
\begin{align*}
\frac{
\vct{u}\cdot\vct{j}\!\left(
\bPsi_{\eqtext{Dirac\_positive\_energy\_axis}}(\pm)
\right)}
{\rho\!\left(
\bPsi_{\eqtext{Dirac\_positive\_energy\_axis}}(\pm)
\right)}
&=\frac{\|\vct{p}\|}{E(\vct{p})},
\\
\frac{
\vct{u}\cdot\vct{j}\!\left(
\bPsi_{\eqtext{Dirac\_negative\_energy\_axis}}(\pm)
\right)}
{\rho\!\left(
\bPsi_{\eqtext{Dirac\_negative\_energy\_axis}}(\pm)
\right)}
&=-\frac{\|\vct{p}\|}{E(\vct{p})},\\
\frac{\vct{j}\!\left(
\bPsi_{\eqtext{Dirac\_positive\_energy\_axis}}(\pm)
\right)}
{\rho\!\left(
\bPsi_{\eqtext{Dirac\_positive\_energy\_axis}}(\pm)
\right)}
&=\frac{\vct{p}}{E(\vct{p})},
\\
\frac{\vct{j}\!\left(
\bPsi_{\eqtext{Dirac\_negative\_energy\_axis}}(\pm)
\right)}
{\rho\!\left(
\bPsi_{\eqtext{Dirac\_negative\_energy\_axis}}(\pm)
\right)}
&=-\frac{\vct{p}}{E(\vct{p})}.
\end{align*}
Thus the component equations and current agree exactly with the intrinsic
projector--ladder derivations above.
